\section{Methods}
\label{sec:methods}

Figure~\ref{fig:framework} shows how candidate generation, mathematical and
numerical checks, measured costs, and routing interact. The checks have
different scopes: a scheme certificate concerns an exact iteration; a
projection contract screens generated code; the final check accepts a point
on the original problem.

\begin{figure}[H]
\centering
\resizebox{\linewidth}{!}{%% Figure 1: OptEvolve overview (TikZ picture only).
%% Compile with framework-standalone.tex, or \input into the paper.
%% Read left to right: problem -> typed choices -> admission -> measurement.
%% The three small cards inside search are the *kinds* of genes, not all eight
%% layers. Figure 2 expands them. Returns below the spine have distinct roles:
%% red = measured kernel cost re-prices the search; dashed = reusable evidence.
\begingroup%
\definecolor{oeInput}{HTML}{E7E0F1}% lavender: supplied problem
\definecolor{oeGenome}{HTML}{DCE6F7}% blue: typed search
\definecolor{oeGate}{HTML}{F4DFA5}% amber: pre-execution admission
\definecolor{oeOutput}{HTML}{E3EFD7}% green: evaluation and routing
\definecolor{oePool}{HTML}{E5E9EA}% gray: retained search evidence
\definecolor{oeInk}{HTML}{293846}% outlines and forward flow
\definecolor{oePrice}{HTML}{B3242B}% device-cost return only
\begin{tikzpicture}[
  x=1cm,y=1cm,font=\sffamily,line cap=round,line join=round,
  panel/.style={draw=oeInk,line width=.8pt,rounded corners=2.4mm,
                align=center,inner sep=1.4mm},
  mini/.style={draw=oeInk!66,line width=.65pt,rounded corners=1.2mm,
               fill=white,align=center,inner sep=1.1mm,
               font=\sffamily\fontsize{7.8}{9.2}\selectfont},
  title/.style={font=\sffamily\fontsize{9.1}{10.8}\bfseries\selectfont,
                align=center},
  note/.style={font=\sffamily\fontsize{7.5}{8.9}\selectfont,
               align=center},
  flow/.style={-{Latex[length=2.3mm]},draw=oeInk,line width=.95pt},
  evidence/.style={-{Latex[length=2.3mm]},draw=oeInk!72,
                    line width=.85pt,dash pattern=on 2mm off 1.2mm},
  price/.style={-{Latex[length=2.3mm]},draw=oePrice,line width=1pt}
]
%% Panel centres and sizes: all four share y=5.50 and height 4.25 cm.
%% Edges are 0.05/2.85, 3.30/6.95, 7.40/10.95, 11.40/14.90:
%% the three forward-arrow gaps are each 0.45 cm.
\node[panel,fill=oeInput,minimum width=2.80cm,minimum height=4.25cm]
  (input) at (1.45,5.50) {};
\node[panel,fill=oeGenome,minimum width=3.65cm,minimum height=4.25cm]
  (genome) at (5.125,5.50) {};
\node[panel,fill=oeGate,minimum width=3.55cm,minimum height=4.25cm]
  (gate) at (9.175,5.50) {};
\node[panel,fill=oeOutput,minimum width=3.50cm,minimum height=4.25cm]
  (evaluate) at (13.15,5.50) {};

%% Input: a family of instances and a deployment target, not a single QP.
\node[title] at (1.45,7.26) {Problem family};
\node[mini,text width=2.32cm,minimum height=.82cm] at (1.45,6.32)
  {Family $P(\theta)$};
\node[mini,text width=2.32cm,minimum height=.82cm] at (1.45,5.28)
  {structure + scale};
\node[mini,text width=2.32cm,minimum height=.82cm] at (1.45,4.24)
  {target device\\accuracy};

%% Typed search: the three decisions that distinguish the proposal from
%% a parameter sweep. The genome schema allows one layer to be patched.
\node[title] at (5.125,7.26) {Typed solver genome};
\node[mini,text width=3.20cm,minimum height=.92cm] at (5.125,6.35)
  {{\bfseries Formulation}\\constraints $\to$ projections};
\node[mini,text width=3.20cm,minimum height=.92cm] at (5.125,5.29)
  {{\bfseries Algorithm}\\scheme + parameters};
\node[mini,text width=3.20cm,minimum height=.92cm] at (5.125,4.23)
  {{\bfseries Execution}\\fused CPU/GPU kernel};

%% Admission has two different checks. The theorem certifies the formal
%% scheme/relaxation grammar; tests against a contract check custom projections.
\node[title] at (9.175,7.26) {Convergence first};
\node[mini,text width=3.12cm,minimum height=1.33cm] at (9.175,6.10)
  {{\bfseries Scheme + relaxation}\\
   symbolic conditions\\$\to$ certificate};
\node[mini,text width=3.12cm,minimum height=1.29cm] at (9.175,4.50)
  {{\bfseries Operator contract}\\
   test custom projections\\reject failures};

%% Evaluation and output: one acceptance check for all candidates, then
%% compare evolved solvers with recognized specialists on measured time.
\node[title] at (13.15,7.26) {Measure and route};
\node[mini,text width=3.07cm,minimum height=1.32cm] at (13.15,6.10)
  {{\bfseries Time to acceptance}\\
   $T=\mathrm{setup}+N\,t$};
\node[mini,text width=3.07cm,minimum height=1.30cm] at (13.15,4.51)
  {{\bfseries Routed solver}\\
   evolved or specialist\\+ routing predicate};

%% Forward spine: the gate precedes compilation and performance evaluation.
\draw[flow] (input.east) -- (genome.west);
\draw[flow] (genome.east) -- (gate.west);
\draw[flow] (gate.east) -- (evaluate.west);

%% Price feedback (red): measured device cost can change which *algorithm*
%% is preferred. This is not just a kernel-tuning step after algorithm search.
\draw[price] (13.15,3.375) -- (13.15,2.84) --
  (5.28,2.84) -- (5.28,3.375);
\node[font=\sffamily\fontsize{7.6}{9}\selectfont,text=oePrice,
      fill=white,inner sep=1pt] at (9.22,2.84)
  {device cost can change the winning algorithm};

%% Evidence feedback (dashed): measured successes and failures persist in
%% the gene pool and inform the next one-layer typed proposal. Keep the pool
%% below the red cost annotation while avoiding a large empty band between them.
\node[panel,fill=oePool,minimum width=8.22cm,minimum height=.91cm]
  (pool) at (7.50,1.85) {};
\node[title] at (7.50,2.03) {Gene pool};
\node[note] at (7.50,1.64) {retain measured successes, failures and provenance};
\draw[evidence] (14.02,3.375) -- (14.02,1.85) -- (pool.east);
\draw[evidence] (pool.west) -- (3.29,1.85) --
  (3.29,3.00) -- (4.10,3.00) -- (4.10,3.375);
\end{tikzpicture}%
\endgroup%
}
\caption{OptEvolve searches typed solver designs for a problem family. The
scheme gate checks convergence conditions before evaluation; generated
projections pass separate contract tests. Measured device costs and outcomes
feed back into search, while accepted solves determine the route.}
\label{fig:framework}
\end{figure}

\subsection{Canonical problem class}
A family supplies matrix-vector and adjoint operators, block metadata, training instances, a target instance and accuracy, and available devices. OptEvolve retains reusable candidates and costs; on the target it returns a route and a point that passes the original-problem check (Section~\ref{sec:exp:protocol}), or reports no answer within budget. Its composite convex model is
\begin{equation}
\min_x f(x)+g(x)+h(Lx),\label{eq:composite}
\end{equation}
where $f$ has a Lipschitz gradient, $g,h$ have computable proximal steps, and 
$L$ is a bounded linear operator with computable adjoint $L^\ast$ and a known norm bound $\ell \ge \|L\|$.
%$L$ has an adjoint and a valid norm bound.
A proximal step minimizes $g(u)+\|u-v\|^2/(2\gamma)$; for a constraint indicator it projects onto the feasible set. An iteration operator $T$ maps the current (possibly primal--dual) state to the next. Splitting constructs $T$ from updates for the separate terms \citep{RyuYin2022LargeScale}.

\subsection{Solvers as typed genomes}
\label{sec:genome}
A genome specifies a base scheme (forward--backward, Douglas--Rachford, Davis--Yin, primal--dual hybrid gradient, or Condat--V\~u), stepsizes, and optional relaxation \citep{ChambollePock2011First,Condat2013Primal,Vu2013Splitting}. Associated genes select a solution-preserving formulation, projection implementation, and backend. Canonical serialization deduplicates proposals; typed mutations cannot silently change the mathematical meaning of a field. The certified grammar is defined in Section~\ref{sec:theory}.

\subsection{Certified method mixing}
\label{sec:mixing}
The gate certifies base schemes and relaxation, not arbitrary blends, adaptive policies, or free-form switching. Conditional blend and safeguarded-step rules appear in Appendix~\ref{app:theory}; they are not used in the reported experiments.

\subsection{The Stage-0 type-check gate}
\label{sec:gate}
Before executing a scheme, the gate checks its stepsize, operator-norm, and relaxation inequalities (Table~\ref{tab:certificates}) and records the resulting exact-iteration construction certificate. Invalid proposals return the violated condition to the proposer rather than being silently clamped. The gate assumes an admissible formulation and exact proximal maps; it is not a proof of floating-point code.

\subsection{Compilation and certificates}
An admitted exact scheme is compiled as a JAX iteration. The construction certificate establishes convergence under Section~\ref{sec:theory}'s assumptions. Runtime residuals and, where available, duality gaps monitor the solve; the reported point must independently pass the original-problem acceptance check. Synthesized numerical projections are run first inside a contract harness, not admitted by the symbolic scheme gate. Appendix~\ref{app:methods} details these distinct checks.

\subsection{Hardware-aware backend genes}
\label{sec:hardware}
The backend genes choose precision, batching, and per-device execution, including a temporally fused GPU tile or an unfused loop. Candidate kernels are checked against reference updates; the exact-operator certificate does not cover their floating-point behavior. Calibration prices each tile on its target device rather than transferring a winner across GPUs (Appendix~\ref{app:methods}).

\subsection{Formulation genes}
\label{sec:formulation}
Formulations can equilibrate rows or move recognized constraints from $h(Lx)$ into the projection in $g$, changing iteration count and cost without changing the underlying solution. Relocation requires a valid recovery map, a new bound for the transformed operator, and an exact mathematical projection; only then can the scheme certificate transfer (Proposition~\ref{prop:formulation-transfer}, Appendix~\ref{app:theory}).

\subsection{Operator synthesis under contract}
\label{sec:synthesis}
When projection time limits a promising route, an LLM proposes a faster implementation for a specified set and device. Before a proposal is used in a solve, a contract harness checks it against a reference projection on tested inputs, including degeneracies. Passing finite tests does not prove equality to the exact projection for all inputs. Admitted implementations are timed on each device (Appendix~\ref{app:methods}).

\subsection{The gene pool, the cost model, and routing}
\label{sec:routing}
An append-only gene pool stores accepted designs, measurements, and rejection feedback. The router estimates $T=\mathrm{setup}+N(\mathrm{formulation})\,t(\mathrm{operator},\mathrm{device})$, using measured iteration counts and calibrated per-device costs, not an operator-norm proxy. It considers evolved candidates and recognized generic or structure-specific engines; unmeasured choices are explored within budget. Every returned answer passes the same original-problem check, regardless of route.

\subsection{Evolutionary loop}
\label{sec:loop}
The LLM proposer receives archived results, a mutation menu, and the genome schema; a random proposer controls for typed-space search at the same evaluation budget. Candidates are admitted before timing, scored by training time to an accepted answer, and selected on held-out instances. Algorithm~\ref{alg:optevolve} ties the search and target routing together; Appendix~\ref{app:methods} gives proposal and fitness details.

\paragraph{Proposal policy and adversarial selection.} Which proposer to ask is
itself a choice, and it can be learned without touching the model. A disjoint
LinUCB contextual bandit \citep{li2010linucb} arbitrates the proposal channel,
deciding at each step whether to query the model for a complete typed design,
mutate one legal gene of the incumbent, or sample a fresh legal design from the
same typed space. It is updated from the executed outcome of each candidate and
leaves the proposer weights frozen, so it changes where the evaluation budget is
spent rather than what the model knows. Selection can also be made adversarial.
The manually authored reference program is frozen as the opening opponent with
its source withheld from the proposer, and a challenger is promoted only if it
certifies every case and then wins a paired contest against the current
champion, the original anchor, and a rotating earlier opponent, under randomized
contest order, with a bootstrap lower bound above $1.05\times$ and no case
regressing by more than $20\%$; the promoted champion becomes the next opponent
and the next parent. On the unmixing class at a fixed budget of eight fresh
evaluations, the bandit channel reached its first certified assembly on the
fourth candidate using two model calls, against the sixth candidate and six
model calls for the frozen loop, so what it buys is queries and evaluations
rather than a better solver. Appendix~\ref{app:llm-loop} reports the
completed fixed-budget controls and the seeds in which uniform
random search also recovers the method.

\begin{algorithm}[H]
\caption{OptEvolve: search, certify, measure, and route.}
\label{alg:optevolve}
\small
\begin{algorithmic}[1]
\Require Problem family, training and target instances, accuracy, devices
\Ensure Reusable gene pool; accepted point and route, or no answer
\While{evaluation budget remains}
  \State Propose a typed scheme, formulation, operator, or backend edit.
  \State Check the exact scheme and formulation; contract-test generated code.
  \If{rejected} \State Record the reason for the next proposal.
  \Else \State Compile, time accepted training solves, and update measured costs.
  \EndIf
\EndWhile
\State Include recognized specialists; explore unmeasured routes within budget.
\State \Return a route and an original-problem-accepted point, or no answer.
\end{algorithmic}
\end{algorithm}

\subsection{Implementation}
JAX compiles the admitted iterations. Search uses Qwen2.5-Coder-1.5B-Instruct for genome proposals and Qwen3.8-27B for projection synthesis; timing is on an RTX~4090 workstation and an H100 node, with CPU specialists run on the corresponding host. Machine, cache, and timing details appear in Appendix~\ref{app:methods} and Section~\ref{sec:exp:protocol}.

% ---------------------------------------------------------------------
% methodology.tex -- the extended genome, the rediscovery argument, and
% the hardware layer.  \input at the end of 3-methods.tex.
% Style rule: never use the em-dash.
% ---------------------------------------------------------------------

\subsection{An exact-path lasso solver as a search-space example}
\label{sec:method:rediscover}

Reachability is not rediscovery: the expert's decisions must be expressible
as mutations, and the fitness must make their effects visible. An exact-path
lasso solver illustrates this design. Each of the four decisions in
Table~\ref{tab:rediscover} becomes a mutation of one layer.

\begin{table}[htbp]\centering\footnotesize
\caption{The four decisive steps of an expert derivation, each reachable
as a single typed mutation. D3 and D4 are visible to any wall-clock
fitness; D5 and D6 are not, which is the reason the fitness changes.}
\label{tab:rediscover}
\begin{tabular}{@{}p{1.42cm}p{1.24cm}p{3.05cm}p{3.30cm}p{3.04cm}@{}}
\toprule
Step & Layer patched & The mutation & Gate rule that keeps it honest &
Why the fitness sees it \\
\midrule
D3 carry the factorization &
3 state &
\texttt{factorization.\allowbreak update}: refactorize to insert\_\allowbreak delete &
drift check on the maintained factor &
per-step cost, visible immediately \\
\addlinespace
D4 free screening bound &
4 restrict &
\texttt{rule}: none to equiangular\_\allowbreak bound &
admitted only if provably safe, or paired with a full re-admission pass,
checked statically &
step-count identity proves exactness, so the win is pure \\
\addlinespace
D5 hybridize &
5 handoff &
\texttt{handoff}: null to \{to: working\_\allowbreak set, predicate\} &
the existing switch\_k rule generalized: exact prefix exempt, tail passes
the full gate &
neither engine dominates, so only a sweep shows it \\
\addlinespace
D6 measure, do not bound &
5 handoff &
\texttt{predicate.\allowbreak source}: bound to counted &
none needed; it is a source swap &
only oracle regret on a sweep sees it \\
\bottomrule
\end{tabular}
\end{table}

The D6 decision illustrates why fitness matters. The gene value
$\texttt{source} \in \{\text{bound}, \text{fitted}, \text{probe},
\text{counted}\}$ is an enumeration, so swapping it is a legal single
mutation. Under a wall-clock fitness, the negated shifted geometric mean
of time-to-certificate over a training set, \texttt{bound} looks fine,
because it is calibrated on the instance where the bound is tight. Under
regret against a per-problem oracle across $54$ shape-and-grid
combinations at two thread counts, \texttt{bound} costs $14\times$ to
$19\times$ and \texttt{counted} costs $1.10\times$ with a worst case of
$1.64\times$. The same mutation gets opposite verdicts under the two
fitness measures; the chosen fitness thus determines which predicate a search
would favor.

The example shows how an expert's four-step derivation can be expressed as
four single mutations; it does not show that a search run rediscovered them.
An ablation of that discovery claim would withhold the four decisive gene-pool
cards, compare LLM mutation without a gene pool, and run a random proposer
over the same typed space. SimpleTES \citep{Ye2026StructuredScaling} is another
algorithm-discovery system with a lasso-path solver and offers an external
comparison target.

\subsection{Where hardware kernel fusion comes in}
\label{sec:method:hardware}

Fusion is layer~8, and layer~8 is not a post-processing step applied to
the winning algorithm. It is one of three hardware genes, all set from the
same quantity, the work per step of whatever inner loop layers~1
through~6 produced.

\begin{table}[htbp]\centering\footnotesize
\caption{The three hardware genes. Each has a GPU form and a CPU form, and
each has a measured cliff, which is what distinguishes a gene from an
environment default.}
\label{tab:hwgenes}
\begin{tabular}{@{}p{1.42cm}p{1.95cm}p{3.35cm}p{2.35cm}p{2.95cm}@{}}
\toprule
Gene & Question & GPU form & CPU form & The cliff that proves it is a gene \\
\midrule
7 admission &
cross the device at all? &
streamed bytes against the bandwidth difference, decided before touching
the runtime &
stay &
importing the runtime costs more than a small solve \\
\addlinespace
8 width &
how much work per launch or barrier? &
fusion tile: $T$ iterations per launch on a $(B{+}2T)^2$ tile with halo in
shared memory &
BLAS thread count from work per step &
the wrong device's tile costs $1.06\times$ to $1.63\times$; $64$ threads
instead of $16$ costs $103\times$ \\
\addlinespace
6 contract &
what precision, and where &
per-expression fastmath scope &
same &
fastmath on a Schur complement turned a $10^{-14}$ certificate into
$10^{-2}$ \\
\bottomrule
\end{tabular}
\end{table}

The fusion tile and the thread count are the same gene at two
granularities. Both answer how wide the parallel region should be relative
to the work the algorithm gives it, both are read off the device profile,
and both have a cliff when set by an environment default instead of by the
algorithm. That is what makes the RTX~4090 against H100 tile result and
the $16$ against $64$ thread result one finding rather than two anecdotes.

Algorithm~\ref{alg:optevolve} closes the loop by feeding the measured
$t(\mathrm{operator},\mathrm{device})$ back into the layer-5 routing
predicate. A better fused tile lowers $t$ for an iterative candidate and
moves the crossover at which it beats another algorithm. Hardware choice
therefore changes the route, not just the cost of a fixed solver
(Equation~\ref{eq:costmodel}).

